\documentclass[a4paper,12pt]{article}

% --- Кодировки и язык ---
\usepackage[T2A]{fontenc}
\usepackage[utf8]{inputenc}
\usepackage[russian]{babel}

% --- Математика ---
\usepackage{amsmath, amssymb, amsthm, mathtools}

% --- Шрифты и типографика ---
\usepackage{helvet}
\renewcommand{\familydefault}{\sfdefault}
\usepackage{icomma}
\usepackage{microtype}
\usepackage{setspace}
\onehalfspacing
\usepackage{parskip}

% --- Графика ---
\usepackage{tikz}

% --- Таблицы ---
\usepackage{tabularx, booktabs, longtable}

% --- Списки и ссылки ---
\usepackage{enumitem}
\usepackage{hyperref}
\hypersetup{
    colorlinks=true,
    linkcolor=darkaccent,
    urlcolor=accent,
    pdftitle={Чек-лист: Тригонометрические уравнения и тождества},
    pdfauthor={Лёвин А.А.}
}

% --- Колонтитулы ---
\usepackage{fancyhdr}
\pagestyle{fancy}
\fancyhf{}
\fancyfoot[C]{\thepage}
\renewcommand{\headrulewidth}{0pt}

% --- Цветовая палитра ---
\usepackage{xcolor}
\definecolor{accent}{HTML}{008080}
\definecolor{darkaccent}{HTML}{004D4D}
\definecolor{textgray}{HTML}{333333}

% --- Поля ---
\usepackage{geometry}
\geometry{a4paper, top=20mm, bottom=20mm, left=20mm, right=20mm}

% --- Настройка заголовков ---
\usepackage{titlesec}
\titleformat{\section}
  {\normalfont\Large\bfseries\color{darkaccent}}
  {\thesection}{1em}{}[\titlerule]
\titleformat{\subsection}
  {\normalfont\large\bfseries\color{accent}}
  {\thesubsection}{1em}{}

\begin{document}

% ==============================================================================
% ТИТУЛЬНЫЙ ЛИСТ
% ==============================================================================
\begin{titlepage}
    \centering
    \vspace*{2cm}
    
    {\Huge\bfseries\color{darkaccent} Чек-лист\par}
    \vspace{0.5cm}
    {\LARGE\color{accent} Тригонометрические уравнения и тождества: \\[0.3cm]
    от суммы кубов до универсальной подстановки\par}
    
    \vfill
    
    {\Large\color{textgray} Лёвин Артём Александрович \\[0.2cm]
    эксклюзивно для Тимофея Васильченко\par}
    
    \vspace{1cm}
    
    {\large\today\par}
    
    \vfill
    
    % Кривая Безье как аккуратный графический элемент
    \begin{tikzpicture}
        \draw[accent, line width=2pt] (0,0) .. controls (4,2) and (8,-2) .. (12,0);
    \end{tikzpicture}
    \vspace*{1cm}
\end{titlepage}

% ==============================================================================
% ОСНОВНОЙ ЧЕК-ЛИСТ
% ==============================================================================
\section*{Стратегия и методы решения}

Данный чек-лист систематизирует продвинутые методы решения тригонометрических уравнений, выходящие за рамки базовых формул. Главная цель --- научиться видеть \textit{структуру} уравнения и выбирать оптимальный путь, избегая стандартных ловушек с ОДЗ.

\subsection*{1. Преобразование разности/суммы в произведение (скрытые приёмы)}
Если в уравнении присутствует выражение вида $\sin x \pm \cos x$, стандартные формулы суммы/разности напрямую не применяются. 
\begin{itemize}[leftmargin=*, itemsep=2pt]
    \item \textbf{Приём:} Используем формулу приведения $\cos x = \sin\left(\frac{\pi}{2} - x\right)$.
    \item \textbf{Алгоритм:} $\sin x - \cos x = \sin x - \sin\left(\frac{\pi}{2} - x\right) = 2\sin\left(x - \frac{\pi}{4}\right)\cos\frac{\pi}{4} = \sqrt{2}\sin\left(x - \frac{\pi}{4}\right)$.
    \item \textbf{Зачем:} Позволяет свести уравнение к виду $\sin A = \sin B$ или разложить на множители.
\end{itemize}

\subsection*{2. Метод вспомогательного угла}
Применяется для уравнений вида $a\sin x + b\cos x = c$.
\begin{itemize}[leftmargin=*, itemsep=2pt]
    \item \textbf{Алгоритм:} Делим обе части на $\sqrt{a^2 + b^2}$.
    \item \textbf{Формула:} $a\sin x + b\cos x = \sqrt{a^2+b^2}\sin(x + \varphi)$, где $\cos\varphi = \frac{a}{\sqrt{a^2+b^2}}$, $\sin\varphi = \frac{b}{\sqrt{a^2+b^2}}$.
    \item \textbf{Пример:} $\sqrt{3}\sin x + \cos x = 2 \implies 2\left(\frac{\sqrt{3}}{2}\sin x + \frac{1}{2}\cos x\right) = 2 \implies \sin\left(x + \frac{\pi}{6}\right) = 1$.
\end{itemize}

\subsection*{3. Универсальная тригонометрическая подстановка}
Позволяет выразить все тригонометрические функции через одну ($t = \tg\frac{x}{2}$). Идеальна для уравнений-«винегретов», где присутствуют $\sin x$, $\cos x$ и $\tg x$.
\begin{itemize}[leftmargin=*, itemsep=2pt]
    \item $\sin x = \frac{2t}{1+t^2}$
    \item $\cos x = \frac{1-t^2}{1+t^2}$
    \item $\tg x = \frac{2t}{1-t^2}$
    \item \textbf{КРИТИЧЕСКАЯ ЛОВУШКА (ОДЗ):} Замена $t = \tg\frac{x}{2}$ сужает область определения! Знаменатель $\cos\frac{x}{2} \neq 0$. 
    \item \textbf{Обязательный шаг:} Решить уравнение $\cos\frac{x}{2} = 0 \implies x = \pi + 2\pi k$. Подставить эти корни в \textit{исходное} уравнение. Если обращает его в верное равенство --- добавить к ответу.
\end{itemize}

\subsection*{4. Формулы сокращённого умножения высшей степени}
Часто встречаются в задачах с кубами тригонометрических функций.
\begin{itemize}[leftmargin=*, itemsep=2pt]
    \item \textbf{Сумма кубов:} $a^3 + b^3 = (a+b)(a^2 - ab + b^2)$.
    \item \textbf{Применение:} $\sin^3 x + \cos^3 x = (\sin x + \cos x)(\sin^2 x - \sin x\cos x + \cos^2 x) = (\sin x + \cos x)(1 - \frac{1}{2}\sin 2x)$.
    \item \textbf{Разность кубов:} $a^3 - b^3 = (a-b)(a^2 + ab + b^2)$.
\end{itemize}

\subsection*{5. Тангенс и котангенс суммы: ОДЗ и ограничения}
Формула: $\tg(\alpha \pm \beta) = \frac{\tg\alpha \pm \tg\beta}{1 \mp \tg\alpha\tg\beta}$.
\begin{itemize}[leftmargin=*, itemsep=2pt]
    \item \textbf{Ловушка 1:} Исходные углы не должны обращаться в $\frac{\pi}{2} + \pi n$. То есть $\cos\alpha \neq 0$ и $\cos\beta \neq 0$.
    \item \textbf{Ловушка 2:} Знаменатель формулы не должен равняться нулю: $1 \mp \tg\alpha\tg\beta \neq 0$.
    \item \textbf{Правило:} Перед решением уравнения с $\tg$ или $\ctg$ \textit{всегда} выписывайте систему ограничений для знаменателей. Найденные корни обязательно проверяйте по этой системе.
\end{itemize}

\section*{Типичные ошибки}
\begin{enumerate}[leftmargin=*, itemsep=2pt]
    \item \textbf{Потеря корней при универсальной подстановке:} Забыли проверить $x = \pi + 2\pi k$.
    \item \textbf{Игнорирование ОДЗ для тангенса:} Получили корень, при котором $\cos x = 0$, и записали его в ответ.
    \item \textbf{Арифметика в формулах суммы/разности:} Ошибки при вычислении $\frac{\alpha \pm \beta}{2}$ и $\frac{\alpha \mp \beta}{2}$.
    \item \textbf{Неправильное раскрытие кубов:} Попытка раскрыть $(a+b)^3$ вместо $a^3+b^3$.
\end{enumerate}

\newpage
% ==============================================================================
% ТРЕНИРОВОЧНЫЙ БЛОК
% ==============================================================================
\section*{Тренировочный блок}

\textit{Решите уравнения, применяя методы из чек-листа. Обратите особое внимание на ОДЗ и проверку корней.}

\vspace{0.5cm}

\textbf{Базовый уровень (закрепление алгоритмов)}
\begin{enumerate}[leftmargin=*, itemsep=8pt]
    \item $\sin x - \cos x = \sqrt{2}\sin 3x$
    \item $\sin x + \sqrt{3}\cos x = 1$
    \item $\sin^3 x + \cos^3 x = 0$
\end{enumerate}

\vspace{0.5cm}

\textbf{Повышенный уровень (комбинация методов и ОДЗ)}
\begin{enumerate}[leftmargin=*, itemsep=8pt, start=4]
    \item $\sin x + \cos x = \sqrt{2}\cos 3x$
    \item $\tg\left(x + \frac{\pi}{4}\right) = 2$
    \item $\sin^3 x - \cos^3 x = \sin x - \cos x$
    \item $\tg\left(x + \frac{\pi}{6}\right) = \sqrt{3}$
    \item $5\sin x + 12\cos x = 13$
    \item $\sin^3 x + \cos^3 x = 1 - \frac{1}{2}\sin 2x$
\end{enumerate}

\vspace{0.5cm}

\textbf{Продвинутый уровень (стратегический выбор)}
\begin{enumerate}[leftmargin=*, itemsep=8pt, start=10]
    \item $\ctg x - \tg x = 2\sqrt{2}\sin 2x$
\end{enumerate}

\newpage
% ==============================================================================
% ОТВЕТЫ
% ==============================================================================
\section*{Ответы к тренировочному блоку}

\begin{table}[h]
    \centering
    \renewcommand{\arraystretch}{1.6}
    \caption*{\large\bfseries\color{darkaccent} Ключи для самопроверки}
    \vspace{0.2cm}
    \begin{tabularx}{\linewidth}{c >{\raggedright\arraybackslash}X}
        \toprule
        \textbf{№} & \textbf{Ответ и комментарии} \\
        \midrule
        1 & $x = -\frac{\pi}{8} + \pi k, \quad x = \frac{5\pi}{16} + \frac{\pi n}{2}, \quad k, n \in \mathbb{Z}$ \\
        \midrule
        2 & $x = -\frac{\pi}{6} + 2\pi k, \quad x = \frac{\pi}{2} + 2\pi k, \quad k \in \mathbb{Z}$ \\
        \midrule
        3 & $x = -\frac{\pi}{4} + \pi k, \quad k \in \mathbb{Z}$ \\
        \midrule
        4 & $x = \frac{\pi}{16} + \frac{\pi k}{2}, \quad x = -\frac{3\pi}{8} + \pi n, \quad k, n \in \mathbb{Z}$ \\
        \midrule
        5 & $x = \arctg 2 - \frac{\pi}{4} + \pi k, \quad k \in \mathbb{Z}$. \textit{ОДЗ проверено.} \\
        \midrule
        6 & $x = \frac{\pi}{4} + \pi k, \quad x = \frac{\pi n}{2}, \quad k, n \in \mathbb{Z}$ \\
        \midrule
        7 & $x = \frac{\pi}{6} + \pi k, \quad k \in \mathbb{Z}$. \textit{ОДЗ проверено.} \\
        \midrule
        8 & $x = \arcsin\frac{5}{13} + 2\pi k, \quad k \in \mathbb{Z}$ \\
        \midrule
        9 & $x = 2\pi k, \quad x = \frac{\pi}{2} + 2\pi k, \quad k \in \mathbb{Z}$ \\
        \midrule
        10 & $x = \pm\frac{\pi}{8} + \pi k, \quad k \in \mathbb{Z}$. \textit{ОДЗ: } $\sin 2x \neq 0 \text{ (выполняется).}$ \\
        \bottomrule
    \end{tabularx}
\end{table}

\vfill
\begin{center}
    \small\color{textgray}
    \textit{Успехов в подготовке! Помните: в тригонометрии нет нерешаемых уравнений, есть не замеченные закономерности.}
\end{center}

\end{document}